\section{Bài tập chuyên đề}

\subsection{Củng cố}

\begin{tuluyen}
Cho các chất sau: \ce{NH4NO3}, \ce{KCl}, \ce{(NH4)2SO4}, \ce{CO(NH2)2}, \ce{K2SO4}. Phân loại các chất thích hợp làm phân đạm hoặc phân kali.
\dapso{Phân đạm: \ce{NH4NO3}, \ce{(NH4)2SO4}, \ce{CO(NH2)2}; phân kali: \ce{KCl}, \ce{K2SO4}.}
\end{tuluyen}

\begin{tuluyen}
Tính khối lượng nitrogen có trong:
\begin{tasks}(2)
\task \SI{50}{\kilo\gram} urea chứa $46\%$ N.
\task \SI{80}{\kilo\gram} phân chứa $21\%$ N.
\task \SI{25}{\kilo\gram} phân chứa $35\%$ N.
\task \SI{120}{\kilo\gram} phân chứa $26\%$ N.
\end{tasks}
\dapso{\SI{23}{\kilo\gram}; \SI{16.8}{\kilo\gram}; \SI{8.75}{\kilo\gram}; \SI{31.2}{\kilo\gram}.}
\end{tuluyen}

\begin{tuluyen}
Một ruộng cần bổ sung \SI{69}{\kilo\gram} nitrogen. Tính khối lượng urea chứa $46\%$ N cần dùng nếu bỏ qua mọi hao hụt.
\dapso{\SI{150}{\kilo\gram}.}
\begin{loigiaiBT}
Khối lượng urea cần dùng:
\[
m=\frac{69}{0.46}=\SI{150}{\kilo\gram}.
\]

\textbf{Vậy,} cần \SI{150}{\kilo\gram} urea.
\end{loigiaiBT}
\end{tuluyen}

\begin{tuluyen}
Một hecta cây cần hấp thụ \SI{72}{\kilo\gram} nitrogen. Dùng phân chứa $30\%$ N, cây sử dụng được $80\%$ lượng nitrogen đã bón. Tính khối lượng phân cần dùng cho \SI{2.5}{\hectare}.
\dapso{\SI{750}{\kilo\gram}.}
\begin{loigiaiBT}
Khối lượng nitrogen cây cần hấp thụ:
\[
m_{N,\text{hấp thụ}}=72\times2.5=\SI{180}{\kilo\gram}.
\]
Khối lượng nitrogen phải bón:
\[
m_{N,\text{bón}}=\frac{180}{0.80}=\SI{225}{\kilo\gram}.
\]
Khối lượng phân:
\[
m_{\text{phân}}=\frac{225}{0.30}=\SI{750}{\kilo\gram}.
\]

\textbf{Vậy,} cần \SI{750}{\kilo\gram} phân.
\end{loigiaiBT}
\end{tuluyen}

\begin{tuluyen}
Viết phương trình hóa học cho các phản ứng sau:
\begin{tasks}(2)
\task \ce{NH4Cl} với \ce{NaOH}.
\task \ce{NH4NO3} với \ce{KOH}.
\task \ce{(NH4)2SO4} với \ce{NaOH}.
\task \ce{NH4Cl} với \ce{Ca(OH)2}.
\end{tasks}
\dapso{Xem các phương trình ở phần lí thuyết.}
\end{tuluyen}

\begin{tuluyen}
Cho \SI{10.7}{\gram} \ce{NH4Cl} tác dụng hoàn toàn với dung dịch \ce{NaOH} dư. Tính thể tích \ce{NH3} thu được ở điều kiện có $V_m=\SI{24.8}{\liter\per\mole}$.
\dapso{\SI{4.96}{\liter}.}
\begin{loigiaiBT}
\[
n_{\ce{NH4Cl}}=\frac{10.7}{53.5}=\SI{0.200}{\mole}.
\]

\[
\begin{array}{c@{\;}c@{\;}c@{\;}c@{\;}c}
\ce{NH4Cl} & + & \ce{NaOH} & \ce{->} & \ce{NaCl + NH3 + H2O}\\
0.200 & & & & 0.200
\end{array}
\]

\[
n_{\ce{NH3}}=\SI{0.200}{\mole}.
\]

\[
V_{\ce{NH3}}=0.200\times24.8=\SI{4.96}{\liter}.
\]

\textbf{Vậy,} thể tích \ce{NH3} thu được là \SI{4.96}{\liter}.
\end{loigiaiBT}
\end{tuluyen}

\begin{tuluyen}
Một mẫu phân ammonium sulfate có khối lượng \SI{13.2}{\gram} chứa tạp chất trơ. Cho mẫu tác dụng hoàn toàn với dung dịch \ce{NaOH} dư, thu được \SI{4.464}{\liter} \ce{NH3} ở điều kiện có $V_m=\SI{24.8}{\liter\per\mole}$. Tính độ tinh khiết của \ce{(NH4)2SO4} trong mẫu.
\dapso{Khoảng $90.0\%$.}
\begin{loigiaiBT}
\[
n_{\ce{NH3}}=\frac{4.464}{24.8}=\SI{0.180}{\mole}.
\]

\[
\begin{array}{c@{\;}c@{\;}c@{\;}c@{\;}c}
\ce{(NH4)2SO4} & + & \ce{2NaOH} & \ce{->} & \ce{Na2SO4 + 2NH3 + 2H2O}\\
x & & & & 2x
\end{array}
\]

\[
2x=0.180\Rightarrow x=\SI{0.090}{\mole}.
\]
\[
m_{\ce{(NH4)2SO4}}=0.090\times132=\SI{11.88}{\gram}.
\]

\[
\text{Độ tinh khiết}=\frac{11.88}{13.2}\times100\%=90.0\%.
\]

\textbf{Vậy,} mẫu phân chứa $90.0\%$ \ce{(NH4)2SO4} theo khối lượng.
\end{loigiaiBT}
\end{tuluyen}

\subsection{Bài tập tổng hợp}

\begin{tuluyen}
Trong nông nghiệp, việc lạm dụng phân bón hóa học trong thời gian dài gây ảnh hưởng xấu đến môi trường, đặc biệt khi sử dụng phân đạm như: \ce{NH4NO3}, \ce{NH4Cl}, \ce{(NH4)2SO4} có thể gây ra ảnh hưởng gì với pH của đất?
\begin{tasks}(2)
\task Làm tăng độ pH của đất, khiến đất chua hơn.
\task Làm tăng độ pH của đất, khiến đất mặn hơn.
\task Làm giảm độ pH của đất, khiến đất chua hơn.
\task Làm giảm độ pH của đất, khiến đất kiềm hơn.
\end{tasks}
\source{Trích nguyên văn Câu 6, đề tuyển sinh chuyên Hóa Lâm Đồng 2026.}
\dapso{Làm giảm pH của đất, khiến đất chua hơn.}
\end{tuluyen}

\begin{tuluyen}
Một trang trại rộng \SI{3.0}{\hectare} cần cây hấp thụ trung bình \SI{64}{\kilo\gram} N trên mỗi hectare. Trang trại dùng urea thương mại chứa $92\%$ \ce{CO(NH2)2}; cây sử dụng được $75\%$ lượng N đã bón. Phân được đóng bao \SI{50}{\kilo\gram}.
\begin{enumerate}
\item Tính khối lượng urea thương mại cần dùng.
\item Tính số bao tối thiểu phải mua.
\end{enumerate}
\dapso{Khoảng \SI{595.7}{\kilo\gram}; cần $12$ bao.}
\begin{loigiaiBT}
Khối lượng N cây cần hấp thụ:
\[
m_{N,\text{hấp thụ}}=3.0\times64=\SI{192}{\kilo\gram}.
\]
Trong \SI{1}{\kilo\gram} phân có lượng N:
\[
0.92\times\frac{28}{60}=0.4293\text{ kg N}.
\]
Lượng N thực tế cây hấp thụ được từ \SI{1}{\kilo\gram} phân:
\[
0.4293\times0.75\approx\SI{0.322}{\kilo\gram}.
\]
Do đó:
\[
m_{\text{phân}}=\frac{192}{0.322}\approx\SI{595.7}{\kilo\gram}.
\]
Số bao:
\[
N=\frac{595.7}{50}\approx11.91.
\]

\textbf{Vậy,} cần khoảng \SI{595.7}{\kilo\gram} phân, tức tối thiểu $12$ bao.
\end{loigiaiBT}
\end{tuluyen}

\begin{tuluyen}
Một mẫu phân chỉ chứa \ce{NH4Cl}, \ce{(NH4)2SO4} và tạp chất trơ. Lấy \SI{20.0}{\gram} mẫu cho tác dụng hoàn toàn với dung dịch \ce{NaOH} dư, thu được \SI{7.44}{\liter} \ce{NH3} ở điều kiện có $V_m=\SI{24.8}{\liter\per\mole}$. Một mẫu khác cùng thành phần, cũng có khối lượng \SI{20.0}{\gram}, cho tác dụng với dung dịch \ce{BaCl2} dư thu được \SI{11.65}{\gram} \ce{BaSO4}. Tính khối lượng mỗi muối ammonium trong mẫu ban đầu. Cho $M_{\ce{NH4Cl}}=53.5$, $M_{\ce{(NH4)2SO4}}=132$, $M_{\ce{BaSO4}}=233$ g/mol.
\dapso{$m_{\ce{NH4Cl}}=\SI{10.7}{\gram}$; $m_{\ce{(NH4)2SO4}}=\SI{6.60}{\gram}$.}
\begin{loigiaiBT}
Gọi số mol \ce{NH4Cl}, \ce{(NH4)2SO4} trong mẫu lần lượt là $x$, $y$.

\[
n_{\ce{NH3}}=\frac{7.44}{24.8}=\SI{0.300}{\mole},
\]
\[
n_{\ce{BaSO4}}=\frac{11.65}{233}=\SI{0.0500}{\mole}.
\]

\[
\begin{array}{c@{\;}c@{\;}c@{\;}c@{\;}c}
\ce{NH4Cl} & + & \ce{NaOH} & \ce{->} & \ce{NaCl + NH3 + H2O}\\
x & & & & x
\end{array}
\]
\[
\begin{array}{c@{\;}c@{\;}c@{\;}c@{\;}c}
\ce{(NH4)2SO4} & + & \ce{2NaOH} & \ce{->} & \ce{Na2SO4 + 2NH3 + 2H2O}\\
y & & & & 2y
\end{array}
\]
\[
\begin{array}{c@{\;}c@{\;}c@{\;}c@{\;}c}
\ce{(NH4)2SO4} & + & \ce{BaCl2} & \ce{->} & \ce{BaSO4 + 2NH4Cl}\\
y & & & & y
\end{array}
\]

\[
\begin{cases}
x+2y=0.300,\\
y=0.0500.
\end{cases}
\Rightarrow
\begin{cases}
x=0.200,\\
y=0.0500.
\end{cases}
\]

\[
m_{\ce{NH4Cl}}=0.200\times53.5=\SI{10.7}{\gram},
\]
\[
m_{\ce{(NH4)2SO4}}=0.0500\times132=\SI{6.60}{\gram}.
\]

\textbf{Vậy,} mẫu ban đầu chứa \SI{10.7}{\gram} \ce{NH4Cl} và \SI{6.60}{\gram} \ce{(NH4)2SO4}.
\end{loigiaiBT}
\end{tuluyen}

\begin{tuluyen}
Một nông dân dự định bón phân ammonium sulfate và vôi tôi cho cùng một thửa đất. Hãy đề xuất cách sử dụng hợp lí và giải thích bằng phương trình hóa học vì sao không nên trộn hai chất rồi bón cùng lúc.
\dapso{Nên bón cách thời điểm; vì \ce{(NH4)2SO4 + Ca(OH)2 -> CaSO4 + 2NH3 ^ + 2H2O}, làm thất thoát nitrogen.}
\begin{loigiaiBT}
Nếu trộn trực tiếp:
\[
\ce{(NH4)2SO4 + Ca(OH)2 -> CaSO4 + 2NH3 ^ + 2H2O}.
\]
Khí \ce{NH3} thoát ra làm mất một phần nitrogen của phân đạm.

\textbf{Vậy,} không nên trộn hai chất rồi bón cùng lúc; cần sử dụng cách thời điểm phù hợp.
\end{loigiaiBT}
\end{tuluyen}
